\section{See What Humans Don't See：AI 如何补足人类盲点}

\subsection{超人类感知不是神话，而是某些任务的自然结果}

在讲完机器如何学会人类已经会做的事之后，讲者把问题推进了一步：\textbf{AI 是否应该只停在模仿人类？} 答案显然是否定的。对很多 fine-grained 任务来说，人类本来就不擅长，或者代价太高；此时机器完全可以成为一种``超人类感知工具''。

鸟类细粒度分类、车型识别、工业缺陷检测、医学影像中的微弱模式识别，都属于这一类。它们的共同点是：类别空间很细、视觉差异很小、人类经验不稳定、规模化分析困难。此时机器的价值不在于像人，而在于比人更稳定、更可重复、更能处理海量数据。

\begin{knowledgebox}{什么叫``超人类''}
这里的 superhuman 并不是说模型在所有视觉任务上全面超过人类，而是指在某些\textbf{定义清晰、标注充分、目标单一}的子任务上，模型可以达到高于普通人的分辨率、速度或一致性。
\end{knowledgebox}

\subsection{AI 还能把图像变成社会科学的测量仪器}

本讲里一个很有启发性的例子，是把街景中的汽车识别结果当成``社会透镜''。讲者提到，研究者可以从城市街景中识别车辆类型，再把车型分布与教育水平、收入、投票倾向、环境差异等宏观变量联系起来，从而用视觉信号研究社会结构。\footnote{根据字幕 00:26:28--00:27:24，讲者回顾了利用街景车辆分布推断教育、收入、投票与环境模式的工作。}

这件事非常重要，因为它重新定义了计算机视觉的用途。视觉系统不再只回答``图里是什么''，而可以回答``一个城市或一个社会系统正在呈现什么统计规律''。单个人类观察者无法手工完成这样的大规模汇总，而 AI 能在海量图像中提取稳定的统计结构。

\begin{figure}[H]
\centering
\begin{tikzpicture}[every node/.style={font=\small, align=center}]
\node[draw, rounded corners, fill=yellow!10, minimum width=2.8cm, minimum height=0.95cm] (image) at (0,0) {街景图像\\cars, blocks, roads};
\node[draw, rounded corners, fill=yellow!18, minimum width=3.0cm, minimum height=0.95cm] (vision) at (4.2,0) {视觉识别\\车型、分布、密度};
\node[draw, rounded corners, fill=yellow!26, minimum width=3.0cm, minimum height=0.95cm] (stats) at (8.6,0) {宏观统计关联\\收入、教育、投票、环境};
\node[draw, rounded corners, fill=yellow!8, minimum width=3.4cm, minimum height=0.95cm] (society) at (4.3,-2.0) {AI 成为社会科学的\\高通量测量工具};
\draw[-{Latex[length=3mm]}, thick] (image) -- (vision);
\draw[-{Latex[length=3mm]}, thick] (vision) -- (stats);
\draw[-{Latex[length=3mm]}, thick] (stats) -- (society);
\end{tikzpicture}
\caption{视觉系统作为``社会测量仪器''：从街景对象识别走向城市级统计分析}
\end{figure}

\begin{importantbox}{这里的关键不是相关性本身，而是尺度}
对单个样本做判断，人类并不弱；但在跨城市、跨区域、跨时间的大规模观察上，AI 可以把原本不可测的社会模式变成可分析数据。这是``看见人类看不见的东西''的一种典型形式。
\end{importantbox}

\subsection{人类视觉其实远没有我们想象得可靠}

讲者随后故意做了几个经典实验，目的不是炫心理学常识，而是提醒我们：\textbf{人类视觉的成功经验容易让人忘记自己的局限。} Stroop test 告诉我们注意力会被自动化阅读过程干扰；change blindness 告诉我们即使场景发生显著变化，人类也可能需要较长时间才察觉。\footnote{根据字幕 00:27:24--00:28:47，讲者先用 Stroop test 展示视觉注意力冲突，再用 engine 消失的 change blindness 实验说明大变化也可能被忽略。}

这些实验之所以重要，是因为它们把``人类有偏差''从抽象判断变成了可观察事实：

\begin{itemize}
    \item 我们的注意力资源有限。
    \item 我们会受先验和自动加工过程影响。
    \item 我们在持续监控任务中容易疲劳。
    \item 我们对显著变化的察觉并不总是及时。
\end{itemize}

\begin{figure}[H]
\centering
\begin{tikzpicture}[every node/.style={font=\small, align=center}]
\node[draw, rounded corners, fill=red!8, minimum width=2.8cm, minimum height=0.9cm] (attn) at (0,0) {注意力有限\\cannot watch everything};
\node[draw, rounded corners, fill=red!14, minimum width=2.8cm, minimum height=0.9cm] (prior) at (4.4,0) {先验偏置\\we see expected patterns};
\node[draw, rounded corners, fill=red!20, minimum width=2.8cm, minimum height=0.9cm] (fatigue) at (8.8,0) {疲劳与一致性问题\\errors accumulate};
\node[draw, rounded corners, fill=red!10, minimum width=4.2cm, minimum height=1.0cm] (risk) at (4.4,-2.0) {高风险环境下，人类视觉局限会转化为\\真实的医疗、安全与公平问题};
\draw[-{Latex[length=3mm]}, thick] (attn) -- (risk);
\draw[-{Latex[length=3mm]}, thick] (prior) -- (risk);
\draw[-{Latex[length=3mm]}, thick] (fatigue) -- (risk);
\end{tikzpicture}
\caption{人类视觉的三类局限：注意力、先验偏置与疲劳一致性问题}
\end{figure}

\begin{warningbox}{``人类在环'' 并不自动等于安全}
很多系统设计里喜欢说``最后让人来确认''，好像只要有人盯着屏幕就能兜底。但如果任务本身就是持续、重复、长时间、高压力的视觉监测，那么人类很可能恰恰是最不稳定的那一环。
\end{warningbox}

\subsection{医疗是 AI 补盲最有说服力的应用场景}

讲者随后把话题从心理学实验迅速转到医疗现场，这个转折非常有力。因为在医疗场景里，注意力缺口和人为遗漏不再只是``有趣现象''，而会变成病人风险、流程延迟与系统成本。字幕里提到，美国医疗系统中的 medical errors 长期是极严重的问题，而手术室里对纱布、针、器械的人工盘点，仍然高度依赖手工 checklist。\footnote{根据字幕 00:28:47--00:30:56，讲者以手术室物品追踪、遗留物风险与人工 checklist 的局限，说明 AI 辅助监测的实际价值。}

这里 AI 的角色不是取代医生，而是把医生、护士和手术团队从低效、重复、易漏的视觉盘点中解放出来。理想的系统会做几件事：

\begin{enumerate}
    \item 持续监测关键物品是否进入和离开工作区。
    \item 在出现数量不匹配时实时预警，而不是收尾时才发现。
    \item 把高频重复工作交给机器，让医护把注意力保留给判断与沟通。
\end{enumerate}

\begin{figure}[H]
\centering
\begin{tikzpicture}[every node/.style={font=\small, align=center}]
\node[draw, rounded corners, fill=green!8, minimum width=2.7cm, minimum height=0.9cm] (sense) at (0,0) {环境感知\\camera / sensor};
\node[draw, rounded corners, fill=green!14, minimum width=2.7cm, minimum height=0.9cm] (track) at (3.6,0) {检测与追踪\\item / action};
\node[draw, rounded corners, fill=green!20, minimum width=2.7cm, minimum height=0.9cm] (alert) at (7.2,0) {风险告警\\missing / anomaly};
\node[draw, rounded corners, fill=green!10, minimum width=3.6cm, minimum height=0.9cm] (human) at (10.9,0) {医护决策\\confirm / intervene};
\draw[-{Latex[length=3mm]}, thick] (sense) -- (track);
\draw[-{Latex[length=3mm]}, thick] (track) -- (alert);
\draw[-{Latex[length=3mm]}, thick] (alert) -- (human);
\end{tikzpicture}
\caption{医疗 ambient intelligence 的基本闭环：持续感知、稳定追踪、异常告警、由人完成最终决策}
\end{figure}

\begin{knowledgebox}{为什么医疗监测特别适合 human-centered AI}
\begin{itemize}
    \item 高风险，但又有大量重复视觉劳动。
    \item 对一致性要求很高，人类疲劳成本很大。
    \item 最终责任必须由专业人员承担，因此天然适合``机器辅助 + 人类决策''。
    \item 价值清晰：减少错误、降低等待、提升照护覆盖面。
\end{itemize}
\end{knowledgebox}

\subsection{Hand Hygiene：为什么这个案例非常典型}

讲者在医疗案例里首先讲了 hand hygiene，而不是更 ``高级'' 的诊断任务，这个选择其实非常说明问题。洗手这件事看起来简单，但对降低 hospital-acquired infection 极其关键；与此同时，它又特别适合暴露人类审计体系的局限。医院当然可以雇审计员去盯，但这会立刻遇到三个问题：人手不够、成本过高、长期观察的一致性极差。\footnote{根据字幕 00:42:11--00:44:36，讲者说明 hospital-acquired infection 的严重性、人工 auditor 的局限、RFID 的非特异性，以及 privacy-preserving depth sensor 的优势。}

讲者还顺手对比了一个常见但不够好的技术方案：RFID。一个佩戴工牌的人靠近洗手池或消毒液分发器，并不等于他真的完成了 hand hygiene 动作。这说明高风险场景里，\textbf{邻近性信号不等于行为信号。} 视觉系统的价值，恰恰在于它能够直接分类 ``是否真的发生了洗手动作''，而不是靠弱代理变量猜测。

\begin{importantbox}{为什么深度传感器比普通摄像头更合适}
\begin{itemize}
    \item 它保留了动作轮廓和空间位置，足以做行为识别。
    \item 它天然弱化面部与身份细节，更容易满足隐私要求。
    \item 在 hospital corridor 这种狭小空间里，它比简单 proximity sensing 更接近真实任务目标。
\end{itemize}
\end{importantbox}

讲者给出的一个很有冲击力的结论是：如果拿同一段视频去做标注，算法的一致性明显高于单个人类观察者，甚至要汇总多个人的判断才能接近 AI 的稳定性。这正是 ``see what humans don't see'' 的现实含义：\textbf{不是 AI 比人更懂医疗，而是 AI 在重复、细粒度、长期视觉审计上比人更稳。}

\subsection{ICU 与 Aging in Place：AI 瞄准的是 healthcare 的 dark spaces}

讲者接着把 ``dark spaces of healthcare'' 这个概念讲得很具体。所谓 dark spaces，不一定是灯光昏暗的地方，而是那些\textbf{持续重要、风险很高、又很难长期获得高质量人工监测}的空间。ICU 是典型例子：病人生死攸关，资源消耗极大，医护人力高度紧张，而很多关键恢复指标又依赖持续观察。\footnote{根据字幕 00:44:48--00:47:30，讲者讨论 ICU mobilization、labor shortage、smart sensor 在 ICU 与老年居家照护中的作用。}

其中一个具体任务是 \textbf{mobilization}，也就是帮助 ICU 病人安全地完成移动和起身。对普通人来说，下床、上床、起身坐下似乎只是日常动作；但对 ICU 患者，这些动作既代表恢复进展，也伴随很高风险。AI 在这里可以监测 ``getting out of bed / getting in bed / getting out of chair / getting in chair'' 等事件，并把这些信息交给医护做更好的恢复管理。

而到了 aging in place 场景，目标又略有不同：不是医院内的高密度照护，而是让老年人能够更安全、更独立地留在家中生活。此时 thermal camera、mobility monitoring、sleep pattern、dietary pattern 等传感器信号，构成的是一种\textbf{早期异常发现系统}。它不一定给出诊断，但可以在系统过载之前发出足够早的风险信号。

\begin{figure}[H]
\centering
\begin{tikzpicture}[every node/.style={font=\small, align=center}]
\node[draw, rounded corners, fill=green!10, minimum width=2.8cm, minimum height=0.9cm] (icu) at (0,0) {ICU\\监测恢复与移动};
\node[draw, rounded corners, fill=green!16, minimum width=2.8cm, minimum height=0.9cm] (ward) at (4.4,0) {病房\\监测 hygiene 与流程风险};
\node[draw, rounded corners, fill=green!22, minimum width=2.8cm, minimum height=0.9cm] (home) at (8.8,0) {家庭\\监测老年人长期安全};
\node[draw, rounded corners, fill=green!8, minimum width=4.4cm, minimum height=1.0cm] (dark) at (4.4,-2.0) {这些场景的共同点：\\重要、长期、稀缺人力、不能靠人工持续盯守};
\draw[-{Latex[length=3mm]}, thick] (icu) -- (dark);
\draw[-{Latex[length=3mm]}, thick] (ward) -- (dark);
\draw[-{Latex[length=3mm]}, thick] (home) -- (dark);
\end{tikzpicture}
\caption{healthcare dark spaces：病房、ICU 与家庭照护都需要长期稳定的感知支持}
\end{figure}

\subsection{偏见不是边角问题，而是视觉系统的结构性风险}

讲者在这一段里把另一类``看不见''也点了出来：\textbf{我们经常看不见自己视觉系统里的偏见。} 棋盘阴影错觉等视觉幻觉说明，人类视觉天然带着关于光照、几何与世界结构的内建假设；而在人脸识别、公平性问题中，数据分布与社会历史又会进一步把偏差放大。\footnote{根据字幕 00:31:01--00:33:00，讲者用 checker shadow illusion 说明人类视觉系统的先验偏差，并转向 face recognition fairness 的社会后果。}

这段话最重要的不是``AI 也有 bias''，而是：\textbf{AI 会放大它所吸收的人类偏差。} 如果训练数据对某些肤色、性别、年龄群体覆盖不足，那么模型在这些人群上的错误率就会更高；当系统被部署到自动驾驶、医疗或公共服务场景时，这种差异会转化为现实世界的不公平。

\begin{importantbox}{Bias 的真正危险}
偏见可怕的地方不在于它让系统``不够完美''，而在于它会把历史中的不平等编码进一个看起来客观、自动、可扩展的系统里。规模化之后，伤害会比单个人类判断更广。
\end{importantbox}

\subsection{隐私保护不是让系统少看一点，而是重新设计它怎么看}

紧接着，讲者提出另一种更微妙的``不看见''：有些时候，\textbf{我们不是想让 AI 看得更多，而是想让它只看见对任务必要的那一部分。} 医疗病房、家庭监护、老人照护、儿童活动识别，都存在这样的张力：你需要模型理解行为，却不希望它暴露身份、面部、家庭内部布局等高度私密信息。\footnote{根据字幕 00:33:00--00:36:58，讲者系统讨论 visual privacy，包括 blurring、masking、dimensionality reduction、federated learning、encryption 以及硬件与软件联合设计的 privacy-preserving action recognition。}

这类问题最容易被误解成``把图像打码就好了''。但讲者特别强调，真正有价值的方案往往要把硬件与软件一体设计：例如专门设计镜头或编码方式，让传感器从源头就只保留动作结构所需的信息，而不保留人脸身份或室内细节。

\begin{figure}[H]
\centering
\begin{tikzpicture}[every node/.style={font=\small, align=center}]
\node[draw, rounded corners, fill=cyan!8, minimum width=2.8cm, minimum height=0.95cm] (raw) at (0,0) {原始视频\\信息最完整};
\node[draw, rounded corners, fill=cyan!14, minimum width=2.8cm, minimum height=0.95cm] (naive) at (3.8,0) {简单模糊/遮挡\\隐私提高但任务性能下降};
\node[draw, rounded corners, fill=cyan!20, minimum width=3.2cm, minimum height=0.95cm] (joint) at (8.4,0) {硬件+软件协同编码\\保留动作，压制身份};
\node[draw, rounded corners, fill=cyan!10, minimum width=3.4cm, minimum height=0.95cm] (goal) at (4.2,-2.0) {目标：看见``行为''，尽量不看见``身份''与``私域细节''};
\draw[-{Latex[length=3mm]}, thick] (raw) -- (naive);
\draw[-{Latex[length=3mm]}, thick] (naive) -- (joint);
\draw[-{Latex[length=3mm]}, thick] (joint) -- (goal);
\end{tikzpicture}
\caption{视觉隐私保护的设计思路：不是简单减信息，而是面向任务重构可见性}
\end{figure}

\begin{knowledgebox}{隐私友好视觉系统的设计原则}
\begin{itemize}
    \item 从源头约束：传感器或编码阶段就控制可见信息。
    \item 只保留任务必要特征：例如动作轨迹，而不是完整身份。
    \item 兼顾可用性：保护隐私不能把任务信号也一起毁掉。
    \item 把 privacy 当成系统目标，而不是部署后的补丁。
\end{itemize}
\end{knowledgebox}

\subsection{这一部分的真正结论}

``See what humans don't see'' 表面上讲的是 superhuman vision，实际上包含了三层不同含义：

\begin{enumerate}
    \item \textbf{比人看得更细}：在 fine-grained 任务上实现高分辨率区分。
    \item \textbf{替人看得更稳}：在长期监测与高风险环境中减少遗漏与疲劳。
    \item \textbf{按人要求地看}：在隐私、公平和责任约束下，只看该看的部分。
\end{enumerate}

\begin{warningbox}{技术能力越强，越不能省掉 human perspective}
因为 AI 不只是一个新传感器，它是一个会被规模化部署的判断系统。它可以帮我们补盲，也可以把偏差、监控和不公平放大。能力和约束必须一起设计。
\end{warningbox}

